\section*{Chapter \ref{ch:iv:the-process-by-firm-type} --- The Process by Firm Type}

\begin{solution}{iq:iv:the-process-by-firm-type:1}
Each stage costs more than the one before: a timed test costs the firm almost nothing per candidate, a screen
one interviewer's half-hour, a final round several senior people's day. With a low base rate of eventual
hires, a cheap test with reasonable validity removes most of the candidates who would fail later, so the
expensive stages are spent on the few who might pass. In the other order, the final-round interviewers would
spend most of their time on candidates a twenty-minute test would have stopped, and the firm would interview
fewer people in depth for the same cost.
\iqlookfor{the funnel as a sequence of measurements ordered by cost, and the role of the base rate.}
\end{solution}

\begin{solution}{iq:iv:the-process-by-firm-type:2}
Spend most of the four weeks on arithmetic. The offer probability is the product of the stage pass rates, so
raising the test's pass rate from 0.3 to 0.6 doubles the chance of an offer, whereas game skill only counts if
the test is passed. Daily timed drills with the methods of \cref{ch:iv:mental-arithmetic} (three weeks),
then one week split between screen questions (probability at speed) and one or two practice games to keep the
strength warm.
\iqlookfor{prioritising the binding stage of the funnel over the candidate's comfort zone.}
\end{solution}

\begin{solution}{iq:iv:the-process-by-firm-type:3}
The intern works on one or two desks or projects with a defined deliverable, is reviewed by the people they
worked with (often in a written mid-summer and end-of-summer review) and presents the work at the end. A
return offer is decided by a committee from those reviews (One Quant Book~17, chapter~28). What is scored is
what a first year needs: reliability (tasks done on time without being chased), speed of learning, judgement
about when to ask, the quality of the final deliverable, and conduct with colleagues at every level.
\iqlookfor{the internship as a ten-week work sample scored by the people who will work with the hire.}
\end{solution}

\begin{solution}{iq:iv:the-process-by-firm-type:4}
Waiting is worth $0.4 \times 130 + 0.6 \times 70 = 94 < 100$: accept, if no extension is possible.
Indifference at $q^\ast = (100 - 70)/(130 - 70) = 0.5$. With a one-week extension the pending process
finishes before the choice, so the candidate takes 130 if it comes and 100 otherwise: $0.4 \times 130 + 0.6
\times 100 = 112$. The extension is worth $112 - 100 = 12$, and asking for it costs a phone call.
\iqlookfor{an expected-value comparison, the break-even probability, and the value of the option to wait.}
\end{solution}

\begin{solution}{iq:iv:the-process-by-firm-type:5}
It verifies that the returns are real and yours: statements or administrator reports for the book you ran,
the capital, the gross and net exposure, the turnover, the costs assumed, and references from your risk
manager and colleagues. It will be sceptical about attribution (a team's returns claimed by one person),
about the period (a record that coincides with one favourable regime), about capacity (a Sharpe ratio earned
on a small book), and about statistical strength: three years of monthly returns at a Sharpe ratio of 1.5 has
a standard error of about 0.6 (\cref{ex:iv:the-process-by-firm-type:record}). Prepare a two-page summary with
monthly returns, drawdowns, exposures and a capacity estimate, and know what you may disclose under your
current employer's confidentiality terms (One Quant Book~16, chapter~11).
\iqlookfor{the elements of track-record due diligence and an honest view of its statistical strength.}
\end{solution}

\begin{solution}{iq:iv:the-process-by-firm-type:6}
Hires per intern are $0.6 \times 0.8 = 0.48$, so $30/0.48 = 62.5$: take 63 interns. The acceptance rate is
usually the cheaper lever: it is lost to competing offers, so make return offers earlier in the season,
match the market's terms and keep in touch with those who have accepted until they start. The return-offer
rate is set by the intake's quality and the firm's bar; raising it by lowering the bar changes the hires,
not only their number.
\iqlookfor{the arithmetic of the conversion chain and the difference between a rate you can buy and one that
measures quality.}
\end{solution}

\begin{solution}{iq:iv:the-process-by-firm-type:7}
Decline the first two offers, then accept the first offer better than both. The chance of taking the best is
$\tfrac{r-1}{n}\sum_{i=r}^{n}\tfrac{1}{i-1}$ with $n = 5$, $r = 3$: $\tfrac25(\tfrac12 + \tfrac13 +
\tfrac14) = \tfrac{13}{30} \approx 0.433$, against $\tfrac15$ for accepting at random; a brute force over all
120 orders confirms it. As $n$ grows the optimal rule declines about $n/e$ (37\%) of the offers and the chance
tends to $1/e \approx 0.368$ (for $n = 100$: decline 37, chance 0.371). The model assumes no recall and only
ranks; with deadlines that can be extended, compare offers directly.
\iqlookfor{the best-choice problem, its threshold rule, the exact small-case answer and the $1/e$ limit.}
\end{solution}

\begin{solution}{iq:iv:the-process-by-firm-type:8}
First the preferred firm, today: tell its recruiter that you hold an offer expiring on Friday and ask
whether the decision after Monday can be made by Thursday. Then the second firm: thank it, say you are in a
final process elsewhere and ask for an extension to the following Wednesday. Then the third firm: schedule
its final round as early as possible next week, and say truthfully that you have a deadline; if the first two
conversations settle the matter, withdraw politely. Never invent an offer or a deadline, and never accept an
offer you intend to renege on (\cref{ch:iv:offers-compensation-and-non-competes}): the firms talk to each
other's former employees, and the industry is small.
\iqlookfor{sequencing by preference, asking for time early, and complete honesty about other processes.}
\end{solution}
